\section{Scope \& Setup}
\label{sec:scope}

This paper uses only existing repository artifacts and does not report any new experiment.
Its evidential core is the Candidate B bundle under
\code{results/20260401/candidate\_b\_minlib\_granularity/},
read together with the companion repository notes in
\code{docs/three\_row\_test\_completed.md},
\code{docs/candidate\_b\_minlib\_granularity.md},
\code{docs/candidate\_a\_collapse\_note.md}, and
\code{RESEARCH\_STATUS.md}.

\paragraph{Base-case scope.}
All theorem statements and empirical claims are restricted to the Sen24 base case \((n,m)=(2,4)\).
The paper does not use neighboring-case evidence as part of the main theorem.
Those neighboring-case results remain relevant only as motivation for the pivot from Candidate A to Candidate B.

\paragraph{Local-rationality scope.}
The levers \code{no\_cycle3} and \code{no\_cycle4} are treated throughout as local forbidden-pattern approximations.
They are not presented as full \code{SocialAcyclic}, and nothing in this paper should be read as a full-acyclicity or full-Sen-family claim.
This scope restriction is part of the theorem statement's intended meaning, not merely a footnote.

\paragraph{Lever notation.}
We write \(\Lambda\) for a finite lever family and \(S \subseteq \Lambda\) for an active lever set.
In the bundled witness family, the relevant levers are \code{asymm}, \code{un}, \code{minlib}, \code{no\_cycle3}, and \code{no\_cycle4}.
In the split witness family, \code{minlib} is replaced by \code{decisive\_voter0} and \code{decisive\_voter1}.

\paragraph{CNF notation.}
For a lever-structured encoding map \(\varphi\), the symbol \(\varphi(S)\) denotes the CNF instance produced from the active lever set \(S\).
We write \(\mathsf{SAT}\) and \(\mathsf{UNSAT}\) for the satisfiability status of that instance.
When \(R \subseteq S\), the expression \(S \setminus R\) denotes removal of the levers in \(R\) before re-encoding.

\paragraph{Repair notation.}
The object of interest is not a clause-level MCS projection.
It is a lever-level repair family associated with an encoded UNSAT instance.
This paper calls that object the raw repair explanation.
Grouped repair explanation is a separate presentation layer that is discussed only after the raw object has been fixed.

\paragraph{Repository setup.}
The manuscript uses the shared repository infrastructure for encoding, auditing, and artifact generation, but its paper files are confined to \code{papers/m1\_5/}.
That separation mirrors the conceptual separation developed here:
the underlying finite witness infrastructure may be shared, while the repair presentation layer requires its own explicit contract.
